\section{Objectifs}

\begin{itemize}
\item Distinguer proportion, pourcentage et évolution.
\item Calculer des évolutions successives et réciproques à l'aide des coefficients multiplicateurs.
\item Interpréter moyenne, médiane, quartiles, écart interquartile et écart type.
\item Exploiter la linéarité de la moyenne et étudier l'effet de l'ajout ou de la suppression d'une valeur.
\item Étudier une série regroupée en classes : histogramme, polygone des fréquences cumulées, moyenne et classe médiane.
\item Comparer deux séries à l'aide de couples d'indicateurs adaptés.
\item Calculer des fréquences marginales et conditionnelles dans un tableau croisé.
\end{itemize}

\section{Prérequis}

\begin{itemize}
\item pourcentages et proportions
\item moyenne, médiane, quartiles
\item lecture de graphiques
\end{itemize}

\section{Activité de découverte}

Deux articles indiquent « +20 \% puis -20 \% ». On vérifie que les deux évolutions ne s'annulent pas. Puis on compare deux séries de données ayant la même moyenne mais des dispersions différentes : un indicateur isolé ne suffit pas toujours à décrire une distribution.

\section{Vocabulaire}

\begin{itemize}
\item \textbf{effectif}, \textbf{fréquence}, \textbf{proportion}
\item \textbf{variation absolue}, \textbf{taux d'évolution}, \textbf{coefficient multiplicateur}
\item \textbf{moyenne}, \textbf{médiane}, \textbf{quartile}, \textbf{écart interquartile}, \textbf{écart type}
\item \textbf{classe}, \textbf{histogramme}, \textbf{fréquence cumulée}
\item \textbf{tableau croisé}, \textbf{fréquence marginale}, \textbf{fréquence conditionnelle}
\end{itemize}

\section{Cours}

\subsection{Proportions et pourcentages}

Une proportion compare une partie à un ensemble de référence. Pour calculer un pourcentage de pourcentage, on multiplie les proportions.

Exemple : 30 \% de 40 \% d'une population représentent \texttt{0,30 × 0,40 = 0,12}, soit 12 \% de la population totale.

\subsection{Évolutions}

Entre une valeur initiale \texttt{V1} et une valeur finale \texttt{V2} :

\begin{itemize}
\item variation absolue : \texttt{V2 - V1} ;
\item coefficient multiplicateur : \texttt{V2 / V1} lorsque \texttt{V1≠0} ;
\item taux d'évolution : \texttt{(V2-V1)/V1}.
\end{itemize}

Des évolutions successives se composent en multipliant leurs coefficients multiplicateurs. L'évolution réciproque utilise l'inverse du coefficient multiplicateur.

\subsection{Statistique à une variable}

La moyenne est sensible aux valeurs extrêmes. La médiane et les quartiles décrivent la position des données ; l'écart type et l'écart interquartile décrivent leur dispersion.

\textbf{Linéarité de la moyenne.} Si toutes les valeurs sont transformées selon \texttt{y = ax+b}, alors la moyenne est transformée de la même manière : \texttt{ȳ = a x̄ + b}.

Ajouter ou supprimer une valeur peut modifier la moyenne et la médiane de manière différente. Il faut donc interpréter les indicateurs en lien avec la série, pas seulement les calculer.

Pour comparer deux séries, on peut associer :

\begin{itemize}
\item moyenne et écart type ;
\item médiane et écart interquartile.
\end{itemize}

\subsection{Séries regroupées en classes}

Pour une série continue regroupée en classes de même amplitude :

\begin{itemize}
\item l'histogramme représente les effectifs ou fréquences par classe ;
\item le polygone des fréquences cumulées permet de lire l'accumulation des observations ;
\item la moyenne peut être calculée à partir de la moyenne de chaque classe et de son effectif ;
\item si seule une répartition uniforme dans chaque classe est supposée, le centre de classe sert de valeur représentative ;
\item la classe médiane est celle dans laquelle l'effectif cumulé franchit 50 \% de l'effectif total ; sous l'hypothèse d'une répartition uniforme, on peut y estimer la médiane.
\end{itemize}

\subsection{Croisement de deux variables qualitatives}

Un tableau croisé répartit les individus selon deux variables qualitatives. Une \textbf{fréquence marginale} concerne une modalité sans condition supplémentaire ; une \textbf{fréquence conditionnelle} est calculée dans une sous-population définie par la condition.

Avant tout calcul, il faut identifier clairement l'ensemble de référence : c'est le dénominateur de la fréquence.

\section{Exemples corrigés}

\begin{itemize}
\item \texttt{+10 \%} puis \texttt{+20 \%} donne \texttt{1,10 × 1,20 = 1,32}, soit une évolution globale de \texttt{+32 \%}.
\item Une hausse de \texttt{25 \%} a pour coefficient \texttt{1,25} ; le coefficient réciproque vaut \texttt{1/1,25=0,8}, soit une baisse de \texttt{20 \%} pour revenir à la valeur initiale.
\item Pour les valeurs \texttt{2,4,6}, la moyenne vaut \texttt{4}. Après transformation \texttt{y=3x+1}, la moyenne devient \texttt{3×4+1=13}.
\end{itemize}

\coursefigure{/images/mathematiques/latex-migration/2nde-statistiques-information-chiffree-figure-1.svg}{Trois classes statistiques de même amplitude avec des effectifs différents.}{Pour des classes de même amplitude, l'histogramme rend directement visibles les différences d'effectifs ou de fréquences.}

\section{Algorithmique en Python}

\subsection{Moyenne pondérée par classes}

\begin{verbatim}
centres = [5, 15, 25]
effectifs = [4, 10, 6]
total = sum(effectifs)
moyenne = sum(c*e for c, e in zip(centres, effectifs)) / total
print(moyenne)
\end{verbatim}

Ce calcul utilise les centres de classes comme valeurs représentatives lorsque l'on adopte l'hypothèse de répartition uniforme dans chaque classe.

\section{Erreurs fréquentes}

\begin{itemize}
\item Additionner des taux successifs au lieu de multiplier les coefficients.
\item Confondre une proportion et un taux d'évolution.
\item Confondre 20 \% et 20 points de pourcentage.
\item Calculer une fréquence conditionnelle avec le total de toute la population au lieu de celui de la sous-population.
\item Comparer deux séries uniquement avec leur moyenne alors que leur dispersion est très différente.
\end{itemize}

\section{Formules et idées à retenir}

\begin{itemize}
\item proportion = partie / ensemble de référence
\item coefficient multiplicateur = valeur finale / valeur initiale
\item taux d'évolution = coefficient multiplicateur - 1
\item évolutions successives : produit des coefficients
\item évolution réciproque : coefficient inverse
\end{itemize}

\section{Synthèse}

Le chapitre relie information chiffrée, statistique descriptive et croisement de données. La priorité est l'interprétation : identifier l'ensemble de référence, choisir les indicateurs pertinents et vérifier que le graphique ou le calcul répond réellement à la question posée.
